\chapter{Implementing Useful Optimizations}
\label{ch:pebble-optimization}

The first optimization should expose weaknesses in the IR interface without requiring a second research project to justify it.
Constant folding and conservative dead-code elimination meet that requirement.
Together they exercise operand inspection, typed replacement, dependency traversal, tombstones, verification, and cross-client equivalence.

\section{Fold only what the semantics permits}
\label{sec:pebble-fold}
The folder handles \code{add}, \code{mul}, and \code{lt} when both operand records are constants.
It computes the result in the byte semantics and keeps the original result type.
It does not infer values through phis, fold branches, exploit algebraic identities with unknown operands, or specialize the function for a test input.
An input that is 9 in the demo remains arbitrary in the compiler contract.

Repeated scans propagate new constants through dependency chains.
For an acyclic expression built in operand-before-user order, one scan can fold every literal-only subexpression.
The implementation still uses a fixed-point loop so the mechanism does not rely on that construction order across different blocks.
Every successful fold changes an arithmetic or comparison opcode into a constant, and no phase changes it back.
There are therefore at most $N$ successful folds.

\begin{algorithm}[htbp]
\caption{Monotone constant folding in the fixed instruction set.}
\label{alg:pebble-fold}
\begin{pseudocode}
\Require verified function with total byte arithmetic
\Ensure all directly constant-foldable supported operations become constants
\Repeat
\State $changed\gets false$
\ForAll{live supported binary operations $v$}
\If{both operands of $v$ are constant records}
\State evaluate the operation with the specified arithmetic
\State replace $v$ with a typed constant, preserving its ID and block
\State clear unused fields; $changed\gets true$
\EndIf
\EndFor
\Until{not $changed$}
\end{pseudocode}
\end{algorithm}

Algorithm~\ref{alg:pebble-fold} has a simple $O(N^2)$ worst-case scan bound, excluding verification.
A use-list worklist can avoid rescanning unrelated instructions, but only after the use lists have a correct update contract.
The bounded implementation favors a scan whose termination argument fits beside its code.

\section{Demand starts at observations and control}
\label{sec:pebble-dce}
The demand roots are every return value and every branch condition.
Following a demanded instruction marks both operands, including both incoming operands of a demanded phi.
The closure stops when a complete scan marks nothing new.
Unmarked values become tombstones.
Because every operation is pure and total, deleting such a definition does not remove an effect or trap.
Keeping branch conditions and all terminators preserves the original execution paths and divergence behavior.

A cycle of phis and arithmetic with no path to a return or branch condition is unmarked and can be removed.
A cycle that determines a loop condition is marked and remains.
This is a dependency argument, not the rule \enquote{delete everything whose result is not returned directly}.
The marking loop performs at most $N$ strict additions to its marked set and has an $O(N^2)$ repeated-scan bound.

\section{Work the transformation}
\label{sec:pebble-opt-worked}
Initially the running function returns value 7, which depends on values 3 and 6.
Value 3 depends on input 1 and constant 2.
Value 6 depends on constants 4 and 5.
All seven values are initially demanded.
Folding changes value 6 to constant 8, removing its operand edges.
The next demand closure marks precisely $\{1,2,3,6,7\}$.
Values 4 and 5 become tombstones.

At input 250, value 3 is still 1 and value 7 is still 8.
At input 9, they are still 16 and 128.
The tests compare all 256 inputs before and after the transformation, then compare the lowered VM results as well.
The local proof covers why each fold is valid; the finite tests check the implementation of that proof on the selected programs.

\section{Toward richer propagation}
\label{sec:pebble-sccp-extension}
Sparse conditional constant propagation would add a value lattice and executable-edge state.
A phi would merge only operands on executable incoming edges, and a discovered constant condition could reduce the edge set.
That design introduces new obligations: unknown input is not an unvisited value, a newly executable edge must revisit dependent phis, and deleting dead blocks must repair every surviving phi.
None of those facts follows from the current literal-only folder.

The next useful pass should be chosen from a client need, not from a list of famous optimizations.
For expression programs, local value numbering may expose more benefit than branch reasoning.
For the loop fixtures, loop-aware constant propagation or identity simplification may be more useful.
Each addition should keep the source oracle and both interpreters in the test loop.

\section{Exercises}
\label{sec:pebble-fourteen-exercises}
\begin{enumerate}
\item \label{ex:pebble-fold-count} Compile and optimize \code{(2+3)*(4+5)}. Predict the live value count and VM instruction count afterward, including the return.
\item \label{ex:pebble-root} Explain why a value used only as a loop's branch condition is live even when the function's return value is constant.
\end{enumerate}

\section{Further reading}
\label{sec:pebble-fourteen-reading}
Kildall gives the foundational dataflow setting for more general propagation~\cite{kildall-1973-dataflow}.
Cousot and Cousot explain why abstract-domain meaning and sound transfer functions matter when extending beyond literal evaluation~\cite{cousot-1977-abstract}.
